\documentclass[12pt]{article}

\usepackage[margin = .8in]{geometry}
\usepackage{amsmath}
\usepackage{graphicx}
\usepackage{multicol, enumerate, tabularx}

\usepackage[parfill]{parskip}

\usepackage{adjustbox, soul}

\usepackage{fancyhdr}
\pagestyle{fancy}

\lhead{\emph{Math F113X Finance 2: APR and Compound Interest}}
\rhead{\emph{lecture notes}}

\renewcommand{\emph}[1]{\textsf{\textbf{#1}}}
\usepackage{tikz}
\usetikzlibrary{calc,trees,positioning,arrows,fit,shapes,through, backgrounds}
\usetikzlibrary{patterns}

\usetikzlibrary{decorations.markings}
\usetikzlibrary{arrows}

\usepackage{pgfplots}

\usepackage{longtable}
\usepackage{tabularx}

\newcommand{\ds}{\displaystyle}
\newcommand{\ans}[1][1in]{\rule{#1}{.5pt}}

\newcommand{\points}[1]{(#1 points.)}		% Trying to be lazy.

\usepackage{array}
\newcolumntype{L}[1]{>{\raggedright\let\newline\\\arraybackslash\hspace{0pt}}m{#1}}
\newcolumntype{C}[1]{>{\centering\let\newline\\\arraybackslash\hspace{0pt}}m{#1}}
\newcolumntype{R}[1]{>{\raggedleft\let\newline\\\arraybackslash\hspace{0pt}}m{#1}}
\newcommand{\red}[1]{\textcolor{red}{#1}}

\newcommand{\be}{\begin{enumerate}}
\newcommand{\ee}{\end{enumerate}}


\begin{document}
\begin{enumerate}
\item  \fbox{\emph{Simple Interest}} Suppose Liz borrows \$1000 and agrees to pay it back in 3 years with 5\% simple interest added each year.
		\be
		\item How much \emph{interest} will she pay? In addition to your answer, write down your computation, including the decimal values you would plug into a calculator.
		\vfill
		\item How much will she owe \emph{in total} at the end of three years? Write down the computation and the answer.
		\vfill
		\ee
\item \fbox{\emph{%Simple Interest with 
Annual Interest Rate (APR)}}  Annual interest rates are pro-rated by time.
	\be
	\item Definition: An APR of r\% paid $n$ times year means 
	
	\vfill
	
	\item So an APR of 5\% paid quarterly means
	
	\vfill
	
%	\item Suppose a savings account advertises an \emph{annual percentage rate} of $6\%$ paid \emph{semi-annually} (two-times-a-year). Suppose the initial deposit (the principal) is \$2000.  For each time period, determine the interest accrued and the total assuming no additional deposits or withdrawals.\\
%	
%	\vfill
%	
%	\vfill
%	
%\begin{tabular}{|l|p{5cm}|p{5cm}|}
%\hline
%Time Period & Interest & Total \\
%\hline 
%6 months&&\\
%\hline
%1 year&&\\
%\hline
%1.5 years&&\\
%\hline 
%2 years&&\\
%\hline 
%3 years (try guessing) &&\\
%\hline
%\end{tabular}
%
%\item In this particular example, how would the total value of the savings account at 3 years be different if the $6\%$ interest was paid annually instead of semi-annually?
%\vfill
	\end{enumerate}
\item Determine $n$ for each word below.\\

\begin{tabular}{c|c|c|c|c|c|c}
word:&annually&semi-annually&quarterly&monthly&weekly&daily\\
\hline
$n$ &&&&&&\\
times per year &&&&&&\\
\end{tabular}


\vspace{.3in}

\fbox{\emph{Compound Interest}} 
	
\item What is \emph{compound} interest and how is it different from \emph{Simple} Interest?
\vfill
\newpage
\item Suppose \$1000 will be invested in an account that accrues \emph{compound} interest at an APR of $3\%$ compounded monthly. We will set up a spreadsheet to model this account assuming no other deposits or withdrawals.
		\begin{enumerate}
		\item Open a new spreadsheet and type column headers in cells \verb 'A1, B1, C1, D1, E1, F1'. (See below.)
		\item In row 2, type the  variables for reference. (hint: \verb'A2' is $r$.)	
		\item In row 3, enter the \emph{start} values for this specific problem. Assume time $t$ is in years.
		\item In rows 4 and 5, enter the values/formulas and the end of 1 month and 2 months.
		\item Use the pull down feature to determine the values in the account at the end of 2 years.
		
	\hspace*{-.6in}\begin{tabular}{c|c|c|c|c|c|c|}
	&A&B&C&D&E&F\\
	\hline
	1&{\tt{Interest Rate}}&{{\tt Times per Year}}&{{\tt Principal}}&{{\tt Time Elapsed}}&{{\tt Interest}}&{{\tt \quad Total \quad}} \\
	\hline
	2&$r$&&&&&\\
	\hline
	3&&&&&&\\
	\hline
	4&&&&&&\\
	\hline
	5&&&&&&\\
	\hline
	\end{tabular}
	\end{enumerate}
\vfill
\item Write out the month-by-month computations and look for a pattern.
\vfill

\vfill

\vfill
\item With a principal of $P$ dollars earning $r\%$ interest compounded $n$ times per year for $t$ years will result in a total of $A$ dollars where

\vfill
\end{enumerate}
\end{document}




 \item \emph{Formula for Compound Interest:} If 
 \begin{itemize}
 \setlength{\itemsep}{0pt} 
   \setlength{\parskip}{0pt} 
\item $P$ represents the \emph{principal} or \emph{present value}, 
\item $r$ represents the \emph{annual interest rate} in decimal form,
\item $t$ represents \emph{time} measured in years,
\item  $n$ represents the \emph{number of compounding periods per year},
  \end{itemize}
   write a formula for the \emph{future value} $A$ under compound interest.
		
	\vfill
	
	\vfill


	\item Suppose a different account offers $6.9\%$ compounded daily with a \$5000 minimum balance. Again assuming a minimum deposit, no additional deposits and no withdrawals, how much would this account have after 10 years? %What do you observe?
	\vfill
	
	
	

	\ee
	
	
\end{enumerate}


\end{document}
